\documentclass{article}
\PassOptionsToPackage{numbers,sort&compress}{natbib}
\usepackage[preprint]{neurips_2026}
\usepackage[T1]{fontenc}
\usepackage[utf8]{inputenc}
\usepackage{amsmath,amssymb}
\usepackage[hidelinks]{hyperref}
% Course identification only; preserve the supplied style's layout and fonts.
\makeatletter
\renewcommand{\@noticestring}{COMS W4776: Neural Networks and Deep Learning. Columbia University. Fall 2026 project proposal.}
\makeatother
\title{ReasonSpace: Extra Tokens and Evidence Sensitivity\\in Compositional Reasoning}
\author{Charles Gao \qquad Tuo Zhang \qquad Jiahao Zhang\\Columbia University}
\begin{document}
\maketitle

\section{Introduction}
Additional token positions may provide a language model with useful computation even when their input symbols convey no task-specific information. However, a higher answer accuracy could also reflect formatting preferences or altered attention rather than more reliable evidence use. We ask: \textbf{Do information-free tokens improve compositional reasoning and evidence sensitivity, and do these effects extend from text to visual representations?} ReasonSpace will test this question with synthetic graph problems whose relevant and irrelevant relationships can be changed independently. Our intended contribution is a controlled analysis of whether filler-induced gains survive these changes, across reasoning depths and matched input representations. We seek informative positive or negative results, not a new claim that filler tokens can help, nor a behavioral proof of an internal reasoning algorithm.

\section{Related Work}
Pause-token training demonstrates that delaying answers can improve performance when models learn to use additional positions \citep{goyal2023}. Pfau et al.\ show that meaningless fillers support computation on specially trained algorithmic tasks, while emphasizing the difficulty of learning this behavior \citep{pfau2024}. Jang et al.\ investigate fillers inserted only at inference and find sensitivity to token identity, count, and placement \citep{jang2025}. More recent synthetic-task experiments include filler effects and mechanistic diagnostics \citep{baherwani2026}. Thus, a punctuation sweep alone would offer limited insight; our focus is the relationship between such interventions and correctness under controlled evidence changes.

Fixed fillers also differ from learned or generated reasoning states. Coconut feeds input-dependent hidden states back as subsequent embeddings \citep{hao2024}, while Mull-Tokens trains multimodal latent tokens \citep{ray2025}. These approaches change training or the computation pathway and are not interchangeable with fixed prompt positions. Ablate-to-Validate shows why accuracy alone can misrepresent the use of continuous thought-token content \citep{zhang2026}. We adopt its caution about content utilization without treating replacement sensitivity as definitive mechanistic evidence. Our text--image comparison will isolate presentation effects on a shared task; it does not establish general visual reasoning or transfer learned through training.

\section{Method and Proposed Investigation}
\textbf{Tasks and counterfactuals.} We will generate 12-node directed graphs with two relation types, each a permutation over nodes. A query specifies a start node and a relation sequence of length 1, 2, or 4; deterministic traversal gives its unique endpoint. Each family contains an original problem, a relevant two-edge target swap that changes the endpoint, and an irrelevant swap that preserves it. Swaps preserve graph size and degrees. Four answer choices are fixed within a family and balanced across families. We will verify all answers, exclude duplicate graph/query structures across development and test splits, and keep related variants together. Neutral labels, shuffled facts, node-renaming controls, and recorded prompt lengths reduce shortcuts and length confounds.

\textbf{Model and interventions.} Our primary candidate is Qwen2.5-VL-3B-Instruct, allowing one open-weight checkpoint for text and vision. The core is inference-only: no fillers versus 16 or 64 dot-token positions after the problem and before the answer cue. Controls use 64 commas, a fixed punctuation mixture, or pre-problem dots. Actual vocabulary IDs and positions will be checked after prompt assembly. A generated intermediate-node sequence provides an explicit-reasoning comparator; its content is input-dependent, unlike fixed fillers. Equal token counts do not imply equal runtime. Forced decoding of the identical filler sequence is only a prefill-equivalence check. Learned embeddings and recurrent latent reasoning are outside the core. Chance, majority-choice, and an exact solver provide sanity baselines.

\textbf{Hypotheses and evaluation.} H1 predicts that 64 post-problem dots improve answer-changing pair correctness over no filler. H2 predicts placement-dependent effects; H3 predicts that gains persist under irrelevant edits and nuisance changes. For original/edited inputs $(x_i,x_i')$ with answers $(y_i,y_i')$, define
\begin{equation}
P_i(c)=\mathbf{1}\!\left[\hat y_c(x_i)=y_i\ \land\ \hat y_c(x_i')=y_i'\right].
\end{equation}
The primary contrast is the mean paired difference $P_i(64)-P_i(0)$, equally weighting depths 2 and 4. We will use fixed-position answer-label logits, audit greedy generation, and report 95\% family-bootstrap intervals. Five percentage points is the prespecified practical-effect threshold; intervals may support, rule out, or leave this benefit unresolved. Secondary metrics include ordinary accuracy, answer-preserving pair correctness, depth effects, intermediate-step validity, latency, and peak memory. Invariance alone can reward constant wrong answers. We will fix prompts, precision, scoring, and model revisions, and correct secondary significance tests for multiple comparisons.

\textbf{Feasibility and extensions.} A pilot of 60 families per depth will validate measurement and identify floor/ceiling effects without selecting for filler gains. Confirmation will use 400--800 families per depth, fixed using pilot-based precision estimates and the compute budget. We provisionally budget 20--40 GPU-hours for text on one 16--24 GB GPU, subject to profiling. The planned vision stage renders equivalent graphs with stable layouts across edits and applies the same fixed intervention. It proceeds after legibility checks, $\geq85\%$ one-hop accuracy, and $\geq90\%$ targeted edge-extraction accuracy on development data. H4 predicts positive visual pair-correctness gains under the unchanged intervention. We will compare paired text/image effects and report perception limitations. Reproducible configurations and predictions will accompany the study. Simulated navigation may later test whether evidence-sensitive gains improve episode success, but is not a committed course deliverable. No preliminary results are claimed.

\bibliographystyle{unsrtnat}
\bibliography{references}
\end{document}
